\section{Criação do Documento de Metadados}

\begin{frame}[fragile]{Sintaxe de Arquivo EFD ICMS IPI}
    \begin{lstlisting}[basicstyle=\ttfamily\footnotesize, frame=single, caption={Exemplo de documento EFD}, label={lst:efd_exemplo}]
|0000|EMP|0|01012024|31012024|EMPRESA TESTE LTDA|12345678901234|12345678901|EM|EMPRESA TESTE |EMPRESA|EMPRESA TESTE LTDA|EMPRESA T|E|1|^CRLF
|0001|0|^CRLF
|0005|EMPRESA TESTE LTDA|EMPRESA |COMERCIAL ABC|EMPRESA TE|FULANO DE TAL|RUA DAS FLORES 123|EMPRESA TES|COMERCIAL A|COMERCIAL ABC|^CRLF
|0100|COMERCIAL ABC|12345678901|EMPRESA TESTE L|12345678901234|COMERCIA|contribuinte|COMERCIAL |12345678901234567890|MARIA SANTOS|FULANO DE T|RUA DAS FLO|FULANO DE TAL|COMERCI|^CRLF
|0990|6|^CRLF
|B001|0|^CRLF
|B020|99|100|SAO PAULO|CO|FU|COM|COMERCIAL|EMPRESA TESTE LTDA|01042024|FULANO|0,00|1,00|99,99|100,00|99,99|100,00|1234,11|0,01|999,99|PRODUTO A1|^CRLF
    \end{lstlisting}
    \fonte{Elaborado pelo autor (2026)}
\end{frame}

\begin{frame}{Sintaxe de Arquivo EFD ICMS IPI}
\textbf{Essa sintaxe possui alguns empecilhos que o presente trabalho necessita lidar:}
    \begin{itemize}
        \item Ausência de uma representação oficial que seja legível por máquina.
        \item Estrutura totalmente implícita.
    \end{itemize}
\textbf{A equipe então decidiu:}
    \begin{itemize}
        \item Fazer a tradução da representação oficial para JSON.
        \item Determinar otimizações prévias que já poderiam ser realizadas para evitar custos de leitura. 
    \end{itemize}
\end{frame}

\begin{frame}[fragile]{JSON de Regras de Metadados}
    \begin{lstlisting}[basicstyle=\ttfamily\footnotesize, frame=single, caption={Modelo JSON de Regras de Metadados para a EFD ICMS IPI}, label={lst:json_metadados_1}]
{
    "nome_tabela": "",
    "registro_principal": "",
    "registros": [{
        "caminho": "",
        "ocorrencia": "",
        "campos": [{
                "nome_campo": "",
                "tipo_dado": ""
        }]
    }],
    ...
    \end{lstlisting}
\end{frame}

\begin{frame}[fragile]{JSON de Regras de Metadados}
    \begin{lstlisting}[basicstyle=\ttfamily\footnotesize, frame=single, label={lst:json_metadados_2}]
    ...
    "colunas": [{
        "nome_coluna": "",
        "tipo_dado": "",
        "origem": "",
        "extras": {
            "caminho": "",
            "nome_campo": ""
        }
    }],
    "chave_primaria": []
}
    \end{lstlisting}
    \fonte{Elaborado pelo autor (2026)}
\end{frame}